\section{Seguridad de la IA y regulación internacional}

\subsection{Dos riesgos centrales}

Hassabis divide el riesgo de seguridad de la IA en dos dimensiones:

\begin{enumerate}
    \item \textbf{Uso malicioso (Misuse)}: personas malintencionadas que utilizan las herramientas de IA para causar daño; esto es, para Hassabis, la \textbf{principal preocupación de seguridad} (primary concern)
    \item \textbf{Alineación técnica (Technical Alignment)}: garantizar que el comportamiento de los sistemas de IA sea congruente con la intención humana
\end{enumerate}

\begin{importantbox}{El uso malicioso es la primera prioridad}
Hassabis coloca el uso malicioso —y no el «descontrol de la IA» abstracto— como la principal preocupación de seguridad. Esto incluye:
\begin{itemize}
    \item El uso de IA para realizar ataques cibernéticos a gran escala
    \item El uso de IA para asistir en el diseño de armas biológicas
    \item El uso de IA para generar contenido de deepfake con fines de ingeniería social
    \item Actores estatales y no estatales que utilizan la IA para potenciar su capacidad de causar daño
\end{itemize}
\textit{Idea clave: la amenaza más urgente no suele ser que la IA misma «se vuelva mala», sino que los seres humanos la utilicen deliberadamente con fines maliciosos.}
\end{importantbox}

\subsection{La certificación «Kite Mark» y la detección de engaño}

Hassabis propone un marco pragmático de verificación de seguridad:

\begin{itemize}
    \item \textbf{Pruebas de referencia sobre conducta engañosa} (benchmarks for deception): detectar de manera sistemática si un modelo de IA posee capacidad o tendencia al engaño
    \item \textbf{Certificación «Kite Mark»}: un sello similar a las marcas de certificación de seguridad de producto; los modelos que superan pruebas independientes obtienen esta marca de certificación de seguridad
    \item \textbf{El papel de verificación de los AI Safety Institutes}: respaldar que los institutos nacionales de seguridad de la IA asuman la función de verificación independiente como terceros
\end{itemize}

\begin{knowledgebox}{¿Qué es «Kite Mark»?}
El Kite Mark es un sello de certificación de seguridad de producto otorgado por la Institución de Normas Británicas (BSI), similar a la certificación CCC de China o al marcado CE de la Unión Europea. Hassabis utiliza esta analogía para señalar que los modelos de IA también necesitan una certificación de seguridad equivalente, otorgada por un organismo independiente: no que el propio desarrollador diga «mi modelo es seguro», sino que, tras someterse a pruebas estandarizadas, un tercero lo respalde.
\end{knowledgebox}

\subsection{La propuesta de un sistema de regulación internacional}

La principal demanda de política de Hassabis es crear un \textbf{organismo internacional de regulación de la IA equivalente al Organismo Internacional de Energía Atómica (IAEA)}.

\begin{importantbox}{La política preferida de Hassabis: una IAEA para la IA}
\textit{La acción de política que Hassabis más desea ver es la creación de un organismo internacional semejante a la IAEA, dedicado a la coordinación y verificación globales de la seguridad de la IA.}

El significado profundo de esta analogía:
\begin{itemize}
    \item \textbf{La lección de las armas nucleares}: la tecnología nuclear fue, en la historia humana, la primera tecnología de «riesgo existencial» que exigió coordinación internacional. La creación de la IAEA tomó más de una década, pero terminó por demostrar su necesidad.
    \item \textbf{La lección de la biotecnología}: la Convención sobre Armas Biológicas (BWC), aunque tiene defectos, ofrece un marco básico de normas internacionales.
    \item \textbf{La particularidad de la IA}: la IA se difunde con mayor facilidad que la tecnología nuclear (no requiere minas de uranio), por lo que la coordinación internacional resulta aún más urgente.
\end{itemize}
\end{importantbox}

\begin{warningbox}{Por qué la autorregulación no basta}
Un argumento común en la industria es «dejemos que nos regulemos a nosotros mismos». Hassabis rechaza explícitamente esta postura:

La seguridad de la IA no puede depender únicamente de la autorregulación de las empresas: se necesitan normas internacionales de carácter obligatorio y un mecanismo de verificación independiente. La razón es sencilla: bajo presión competitiva, un compromiso de autorregulación sin restricciones externas es frágil. Es como si ninguna planta nuclear dijera «no necesitamos las inspecciones de la IAEA, nosotros mismos podemos garantizar la seguridad».
\end{warningbox}

\subsection{Resumen de la sección}

La principal amenaza a la seguridad de la IA es el uso malicioso, y para generar confianza se requieren pruebas de referencia de detección de engaño y certificaciones al estilo «Kite Mark». La demanda de política más central de Hassabis es crear una IAEA para la IA: un organismo de regulación y verificación de la seguridad con fuerza vinculante internacional.

%% ==================== Section 8 ====================
